\documentclass{article}
\usepackage[utf8]{inputenc}
\usepackage{amsmath}
\usepackage{tikz}
\usepackage{pgf}

\usepackage{graphicx,calc}
\newlength\myheight
\newlength\mydepth
\settototalheight\myheight{Xygp}
\settodepth\mydepth{Xygp}
\setlength\fboxsep{0pt}
\newcommand*\inlinegraphics[1]{%
  \settototalheight\myheight{Xygp}%
  \settodepth\mydepth{Xygp}%
  \raisebox{-\mydepth}{\includegraphics[height=\myheight]{#1}}%
}
\newcommand{\coffee}[0]{\inlinegraphics{coffee.png}}
\newcommand{\textcoffee}[0]{\inlinegraphics{text_coffee.png}}
\newcommand{\tea}[0]{\inlinegraphics{tea.png}}

\renewcommand{\figurename}{Figur.}

\usepackage{todonotes}
\usepackage{xspace}
\newcommand{\madsfoek}[0]{Mads Føk\xspace}


\title{Kaffepausen: en Sidste Analyse}
\author{}
\date{}

\begin{document}

\maketitle

I må kende til det vidt spredte udtryk: ``Hvis man har en større mængde data, så skal man skrive en \madsfoek\ artikel''. Nu har jeg efterhånden en hel del data, så det vil jeg naturligvis efterleve.

Her er vi igen, en kaffeartikel mere min ven. Og til alle jer nye læsere som har samlet et farverigt \madsfoek blad op for første gang, så velkommen til!
Denne artikel er en fortsættelse på den efterhånden lange serie af analyser på en gruppes kaffeforbrug på datalogi.
Jeg har prøvet på at lade denne artikel stå så meget for sig selv som den kan, men jeg kan ikke udelukke, at der sniger sig noget intertekstualitet ind fra de tidligere artikler og analyser. Hvis du sidder nu og tænker ``oh nej dog, jeg skal jo finde alle de interne jokes, jeg må hellere læse de gamle udgaver først!'', må jeg nok skuffe dig med at de nok er svære at støve op i printet udgave. Men så er det jo godt at året har slået 2025, og internettet for længst har sin storhed, så jeg kan pege dig i retningen af ``https://www.madsfoek.dk/udgivelser'', hvor alle udgivelserne af \madsfoek er til at finde i digital format. Her kan man både finde de gamle artikler, så man kan læse hele triologien i fire dele, men også en stor bunke gamle Hetz hvis man mere er til det. Her ligger de andre artikler i \madsfoek\ nr. 1 årgang 51 og nr. 1 og 2 i årgang 52.
Men nok om det. Du sidder jo med bladet lige nu, og skal ikke tænke på alt muligt gammelt tekst, men I stedet på den skønne tekst du har lige foran dig! Så læn dig godt tilbage, grib din kaffekop, og lad mig starte min historie, som den jo starter bedst.

På datalogi befinder der sig et magisk rum, hvor de ansatte har muligheden for at hente den skønneste livseliksir: kaffen. Af de to legendariske maskiner løber et overflødighedshorn af kaffe, cappucino, latte machiato, kakao, og andre skønne drikkevarer.

Dengang jeg var yngre, og blot instruktor, sad vi tit i Regnecentralen og arbejdede, og ventede på den næste gang nogen sagde de smukke ord: ``Skal vi ikke lige tage et coffee run?''. Og så kunne man få sig en velfortjent pause, mens man traskede afsted, snakkede og fik noget mere dejlig kaffe. Men ak! Den tid kunne ikke blive ved. Man blev mere ansat, fik kontor, og pludseligt var de daglige coffee runs forsvundet, da alle man plejede at gå afsted med, også sad på hver deres kontor rundt omkring. Man gik da stadig i kaffestuen -- man er jo ikke psykopat -- men at tage derned alene havde ikke den samme charme. Man fik ikke sludret! Det var en mørk tid.

Men ud af mørket fødes lyset: '\coffee\ ping'-chatten.
En chat, hvor man kan sende en \coffee\ emoji, og gennem dette batsignal, fortælle at det er tid til et coffee run. Alle, der gerne ville med, kan så sende en \coffee\ som respons, og på den måde ved man hvem man skal vente på, så man kan nå at møde alle, som skulle med på coffee runnet. Et skønt koncept, der fjernede mørket, og kun lod den mørke kaffe blive tilbage.

Men efter en rum tid, så bliver det jo til nogle pings. Og med mange pings kommer der meget data, som man kan analysere på. For hvornår går folk i kaffestuen? Hvor mange er man afsted af gangen?

Citation på introduktionen: \madsfoek\ nr. 1 årgang 52 artikel ``Kaffepausen: En dybtegående analyse'', som citerer \madsfoek\ nr. 1 årgang 51 artikel ``Kaffepausen: en Analyse''.
Du troede nok jeg ville hoppe i plagietfælden hva? Nej nej, ikke sidste gang og heller ikke denne gang! Der er skam kildehenvisning på!


\section*{Metadiskussion af dataen}

Før man kan lave en analyse på data, skal man nu først have fat i dataen.
Jeg gennemgik i detaljer i den første udgave af denne serie for to år siden hvordan man får fat i sin besked historik fra Facebook som en JSON fil. Denne guide galt stadig sidte år nogenlunde, så vi kunne undgå at diskutere det yderligere der. Men hvis der er en ting, der er konstant på intetnettet, så er det, at man konstant ændrer layoutet på hjemmesider uden nogen særlig grund. Guiden gælder derfor, ak og ve, ikke mere. Dog skal man ikke mange klik ind i indstillingerne, før man rammer en ny og meget fancy søgefunktion, som på ingen tid finder frem til den download funktion, man leder efter. Jeg må, modvilligt, klappe af den nye brugergrænseflade.

I sidste ende får man så pakket det data ned man er interaseret i, sammen med en masse anden data man ikke er interaseret i, og efter et par timer er det hele klar til download. Her får man et par fine zip filer, som kun kræver ganske lidt graveri, før man får fat i beskederne fra \coffee\ ping chatten enkodet som en JSON fil.
Denne JSON fil indeholder en smule metadata for chatten samt en lang liste ``messages'', som indeholder alle beskederne i chatten. Hver af disse indeholder ``sender\_name'' (navnet på hvem der har sendt beskeden) og ``timestamp\_ms'' (unix time i millisekunder). Tidligere blev det diskuteret, hvorfor unix time er absolut overlegen over almindeligt datoformat, og denne holdning holder jeg fast i, givet den stadig er 100\% sand.
Og bare rolig, bare fordi almindeligt datoformat er skrald for computere, har jeg naturligvis stadig taget højde for sommer- og vintertid i oversættelsen til det almindelige datoformat, hvor vi dødelige kan følge med. Eller, det har det indbyggede \texttt{datetime} biblotek i Python taget højde for.
Tekstbeskeder indeholder yderligere feltet ``content'', som er indeholdet af, hvad der er blevet sendt. Det er her, kære læser, at vores kaffe pings gemmer sig. Ja, man skulle jo nok tro, at alt der var i \coffee\ ping chatten, er \coffee\ pings, men her tager man fejl, da det også indeholder en utrolig masse spam. Vi skal derfor filtrere i beskederne, for at få fat i de faktiske \coffee\ pings.

Her er JSON et tekstformat, og da vores \coffee\ pings er emojies, som jo egentligt bare er små billeder, kan de derfor ikke være direkte i JSON filen. Heldigvis er både tekst og billeder for en computer bare bits, og nogle kloge hoveder lavede for lang tid siden UTF-8 formatet, som tillader at enkode emojies effiktivt.
Ud over at understøtte emojies, understøtter UTF-8 formatet også (næsten) alle tegn i verden. I JSON filen er UTF-8 karakterer enkodet med '\textbackslash u00YY', hvor 'YY' er hex enkodningen af den byte der ligger underliggende. Da man kun bruger de nederste to hex værdier kunne man enkode den samme information med '\textbackslash xYY', hvilket Facebook af uransagelige årsager ikke har valgt at gøre.
Dog kan vi nu fint genkende UTF-8 karaktererne i JSON filen. Her sætter man flere af disse hex bytes sammen i træk, da UTF-8 er variablet enkodet, for at lave en emoji. Her ved vi, at den enkodning vi leder efter, er
\begin{center}
	\textbackslash u00e2\textbackslash u0098\textbackslash u0095: \coffee \; .
\end{center}
Dog har størstedelen af de beskeder som enkoder et \coffee\ ping yderligere tilføjet '\textbackslash u00ef\textbackslash u00b8\textbackslash u008f' bag på denne enkodning.
Der er dog ingen synlig forskel på '\textbackslash u00e2\textbackslash u0098\textbackslash u0095' og '\textbackslash u00e2\textbackslash u0098\textbackslash u0095 \textbackslash u00ef\textbackslash u00b8\textbackslash u008f', hvilket i de tidligere udgaver har givet en masse hovedpine for at finde ud af hvad pokker denne ekstra enkodning gør.
Her endte den endelige løsning at gå tilbage til '\textbackslash xYY' formattet for enkodningen, da ingen vidst rigtigt bruger det ueffiktive '\textbackslash u00YY' format. Kimse kimse Facebook.
Men med denne anden enkodning bliver den ekstra enkodning til '\textbackslash xef\textbackslash xb8\textbackslash x8f', og her kunne jeg endeligt med en hurtig Google søgning finde frem til en ordentligt tabel over alle UTF-8 enkodninger med klare beskrivelser på en tysk hjemmeside\footnote{https://www.utf8-chartable.de/unicode-utf8-table.pl?start=65024\&utf8=string-literal}.
Og her fandt jeg at den ekstra enkodning betyder \textasteriskcentered\! indsæt selv passende trommehvirvel\! \textasteriskcentered\
\begin{center}
	\textbackslash u00ef\textbackslash u00b8\textbackslash u008f: VARIATION SELECTOR-16 \; .
\end{center}
Her kan man så bruge hele verdens bedste leksikon Wikipedia\footnote{https://en.wikipedia.org/wiki/Variation\_Selectors\_(Unicode\_block)} til at finde ud af hvad en VARIATION SELECTOR gør.
Her ses det, at det blandt andet bruges til kompatibilitet med eksempelvis sjældne Kinesiske tegn, som blev indgivet til Unicode i år 1992-1998, Myanmarske tegn, og Egyptiske hieroglyffer. Og hertil har de moderne hieroglyffer, nemlig emojis, også en varians selector. VARIATION SELECTOR-16 betyder at emojien skal fremvises som en emoji.
Facebook har derved ikke bare brugt den lange verion af UTF-8 enkodning, men de har også brugt yderligere 3 UTF-8 karakterer på alle vores \coffee\ \emph{emoji} pings, på at sige de skal vises som \emph{emojis}.

Men, som du jo nok tænker kære læser, hvis det er variation 16, hvad gør alle de andre variationer så? Jeg er jo akademisk anlagt, så jeg gik dem naturligvis alle igennem, og fandt frem til at alle VARIATION SELECTORs gør intet ved vores kære \coffee\ emoji, pånær variation 15, som giver en særlig tekst variation, som ser ud som følgende
\begin{center}
	\textbackslash u00e2\textbackslash u0098\textbackslash u0095\textbackslash u00ef\textbackslash u00b8\textbackslash u008e: \textcoffee \; .
\end{center}
Hvis man en dag keder sig gevaldigt meget, kan jeg anbefale at lege lidt med en UTF-8 decoder\footnote{https://mothereff.in/utf-8}, for at se hvordan tekst udgaven af dine yndligs emojies ser ud.

Jeg diskuterede i sidste udgave en mulig mængde af fejlkilder som kan være i \coffee\ ping dataen, da det jo er meget akademisk. Herunder fik jeg nævnt at den største fejlkilde nok er Steffan, da han grundet sit brug af en særlig udgave af Messenger havde sendt 151 \coffee\ pings uden VARIATION SELECTOR-16, som derved var gået tabt.
Hertil vil jeg gerne revurdere denne analyse, da efter sidste artikkel er udkommet, er der nærmest gået sport i \coffee\ ping chatten om at være en større fejlkilde, og jeg kan derfor ikke længere alene skylde skylden på Steffan.
Dog kan jeg opsumere for dig kære læser, at Steffan siden sidst har pinget hele 5 gange med VARIATION SELECTOR-15, blot fordi han kan. De er, naturligvis taget højde for i mit script, og tæller derfor med i den kommende analyse.

Men, netop fordi Steffan tidligere havde pinget med andre enkodninger af \coffee\ emojien, lavede jeg naturligvis en analyse af de pings som faldt udenfor den normale enkodning. Her fandt jeg frem til at der pludseligt nu var 214 underlige \coffee\ pings, hvilket er en pludselig stigning, som ikke helt passede overens med de kun 5 nye \textcoffee\ pings.
Jeg tjekkede derfor afsenderen af disse nye sære \coffee\ pings, og fandt til min skræk og forbløffelse ud af, at jeg selv havde sendt flere skæve \coffee\ pings! Dette kan naturligvis ikke være med vilje, givet jeg aldrig frivilligt ville blive en fejlkilde for mig selv. Dog kan jeg gemme mig i mængden, da næsten alle i chatten havde sendt syrlige \coffee\ pings. Jeg gik derfor skridtet videre, og tjekkede dato og klokkeslet for alle disse mærkværdige \coffee\ pings, og kunne se at de nye løjerlige \coffee\ pings alle var fra meget sent i data perioden. Ved at matche disse datoer op mod den rå JSON data fandt jeg frem til årsagen bag disse unaturlige \coffee\ pings: d. 11. september 2025 kl 12:10:27 blev chat emojien skiftet fra \coffee\ til \coffee\! !
Her kan den opmærksomme læser naturligt se at den gamle \coffee\ emoji indeholder VARIATION SELECTOR-16, mens den nye \coffee\ emoji ikke gør. Og alle pings siden denne opdatering indeholder ingen VARIATION SELECTOR, hvilket præcis beskriver alle de excentriske nye \coffee\ pings.
Det forklarer også hvorfor et ekperiment jeg tidligere lavede for at finde ud af om ekstra enkodningen var forskellen på en almindelig emoji og chat emojien slog fejl: Facebook bruger ikke længere VARIATION SELECTOR-16 på chat emojien (fornuftigt). Dette er så årsagen til at mine nye chats lavet til eksperimentet ikke kunne finde den ekstra enkodning. Men da \coffee\ ping chatten blev kreeret for længe siden hvor Facebook brugte VARIATION SELECTOR-16, har den derfor kunne overleve så længe. Forskning.

Nu fik jeg nævnt at flere har forsøgt at være en fejlkilde i gruppen. Her vil jeg fremhæve det hæderlige forsøg lavet af Lasse. Han prøvede at komme ind på ranglisten, ved blandt andet at pinge med en \emph{te} kop frem for en \emph{kaffe} kop, da han kun går i kaffestuen efter te. Disse emojies ser ud som følgende
\begin{center}
  \textbackslash u00F0\textbackslash u009F\textbackslash u008D\textbackslash u00B5: \tea \; .
\end{center}
Hertil vil jeg pointere at emojies ikke renderes ens på alle enheder, og dette har Facebook udnyttet, ved at lave deres eget emoji sæt, hvor te koppen ser absolut rædsom ud.
Lasse stoppede dog med denne type ping efter kun 3 \tea\ ping, da jeg forklarede ham, at de ikke ville tælle med i analysen i denne udgave. Og nu vi taler om det, skal vi så ikke bare se at komme videre til analysen?


\section*{Dataanalysen}

\coffee\ ping chatten blev født d. 19 september 2022, og tidligere har vi ladet dataperioden løbe i et år af gangen, hvilket for første periode løber fra fødselsdagen frem til og med 18. september 2023. Ved at bruge de samme datoer, men med forskudte årstal, opnåes de tre dataperioder.
Heldigvis er formattet for dataen den samme, og jeg kunne derfor simpelt starte scriptet op fra de forrige år og få den samme analyse frem for den nye data. Dog må jeg indrømme at scriptet efterhånden havde udviklet sig til så meget spagetti, at man kunne diskutere om det var en hidtid uopdaget italiensk grund ret. Jeg valgte derfor at kaste noeget energi efter at gøre scriptet pænere og mere strømlignet, hvilket endte ud med, efter en smule arbejde, at kunne gøre præcis det samme som før. Min indre datalog er glad.

Lad os starte ud med at svare på det mest oplagte spørgsmål: hvor mange \coffee\ pings har der så været det sidste år? 3236. Det giver et gennemnit på ca. 8,87 \coffee\ pings om dagen. Det er en stigning på hele ca. 11\% sammenlignet med forrige år. Til sammenligning var stigningen fra første år til andet år på ca. 6\%. Vi er derfor begyndt at pinge markant mere end tidligere, faktisk så meget, at der nu er i gennemsnit næsten et helt ping mere om dagen i gennemsnit. Ved at prikke rundt i dataen kan det ses, at det som forventet skyldes, at en lille del af gruppen står for langt de fleste \coffee\ pings, men blandt alle disse er der en markant forøgelse.
For overskuelighedens skyld er her antal pings fordelt over de tre dataperioder.
\begin{center}
	\begin{tabular}{l|c|c}
		Periode & \coffee\ pings & \coffee\ \!/dag \\ \hline
		2022 - 2023 & 2751 & 7,54 \\
		2023 - 2024 & 2916 & 7,97 \\ 
    2024 - 2025 & 3236 & 8,87 \\ \hline
		2022 - 2025 & 8903 & 8,12
	\end{tabular}
\end{center}
Med hele 8903 \coffee\ pings over de sidste tre år kan man roligt sige at \coffee\ ping chatten til tider føles som absolut spam. Hertil skal man huske at tillægge at \coffee\ ping chatten også indeholder yderligere beskeder, i et omfang hvor kun 2/3 af beskederne er reele \coffee\ pings.

Men vi er jo ikke kommet her for blot at vide et tal så simpelt som et gement genemsnit! Nej, vi er kommet her for en rigtig analyse over hvornår og med hvem folket går i kaffestuen.
Lad os starte med at tage et kig på fordelingen af \coffee\ pings fordelt på ugedagene hen over de sidste tre år. Denne fordeling kan ses på Figur~\ref{fig:weekday_analysis_multi}, hvor de tre år er vist side om side.
Her kan vi se at alle dage, pånær mandag, er gået op i pings. Dette er dog allerede bekrevet blot fra det
rå antal flere \coffee\ pings, der er i det sidste år.
Tager man i stedet et kig på den relative fordeling, kan vi se at tirsdag er blevet den helt store \coffee\ ping dag med synligt flere pings end resten af ugens dage, i modsætning til forrige år, hvor mandag og tirsdag har været lige store \coffee\ ping dage.
Det ses også at weekenderne har markant færre \coffee\ pings end resten af ugens dage, hvilket er meget betrykkende, da det må betyde, at vi holder bare lidt fri i weekenden.
Fredag er igen en gang den hverdag med færrest \coffee\ pings, hvilket kan tilskrives Fredagscaféen\footnote{I folkemunde kendt som DatBar.}.

% \begin{figure}[h!]
% 	\centering
% 	\resizebox{\textwidth}{!}{\input{plot_weekday_analysis_2022-2025.pgf}}
% 	\vspace{-25pt}
% 	\caption{\protect\coffee\ pings fordelt på ugedag i 2022-2025 perioden.}
% 	\label{fig:weekday_analysis_2022-2025}
% \end{figure}

% \begin{figure}[h!]
% 	\centering
% 	\resizebox{\textwidth}{!}{\input{plot_weekday_analysis_2024-2025.pgf}}
% 	\vspace{-25pt}
% 	\caption{\protect\coffee\ pings fordelt på ugedag i 2024-2025 perioden.}
% 	\label{fig:weekday_analysis_2024-2025}
% \end{figure}

\begin{figure}[h!]
	\centering
	\resizebox{\textwidth}{!}{\input{plot_weekday_analysis_multi_2022-2023_2023-2024_2024-2025.pgf}}
	\vspace{-25pt}
	\caption{\protect\coffee\ pings fordelt på ugedag i alle perioder. Her er den første kolonne for hver ugedag den første periode som var i 2022-2023, og så fremdeles for de to næste kolonner.}
	\label{fig:weekday_analysis_multi}
\end{figure}

At det er Fredagscaféen som kan tilskrives de færre \coffee\ pings, kan bedre ses hvis man tager et skridt længere ind i tiden, og kigger på antal \coffee\ pings fordelt på hver time i løbet af ugen. Denne fordeling kan ses på Figur~\ref{fig:weekday_hour_analysis_2022-2025}, for alle totale \coffee\ pings over de sidste tre år.
Dog ses det også at der et lille bump i \coffee\ pings fredag aften, som initielt kan tilskrives at man til tider fristes til en kop kakao efter et par øl i fredagsbaren. På dette plot er bumpet dog blevet større end tidligere, hvilket nok skyldes at jeg tidligere har pointeret i disse artikler hvor ofte vi går i kaffestuen fredag under baren, hvilket har skapt en selvsuplerende spiral.

Hvis vi dog kigger væk fra \coffee\ pings om natten, men i stedet fokusere på den faktisk arbejdstid, kan det ses at der for hver hverdag er to meget tydeligere søjler. Den første skyldes at vi typisk går i kaffestuen når vi først møder ind, hvilket vi efterhånden gør ganske konsistent. Den anden skyldes at vi går i kaffetuen efter frokost, hvilket vi også gør på nogenlunde samme tid hver dag. Bredten af søjlerne skyldes at vi natuligvis går i kaffestuen på andre tidspunkter end dette, men disse er mindre fastlagte, hvor hullet mellem søjlerne skyldes den førnævnte frokostpause.

\begin{figure}[h!]
	\centering
	\resizebox{\textwidth}{!}{\input{plot_weekday_by_hour_analysis_2022-2025.pgf}}
	\vspace{-25pt}
	\caption{\protect\coffee\ pings fordelt på klokkeslet og ugedag over alle tre år.}
	\label{fig:weekday_hour_analysis_2022-2025}
\end{figure}

Fordelingen ændrer sig en smule hvis man kun kigger på det sidste år\footnote{Dette plot er udeladt, da det ikke har meget nyt at vise sammenlignet med fordelingen for alle tre år.}. Her er de karaktariske to søjler stort set væk for tirsdag og til dels torsdag. Dette skyldes med høj sandsynlighed undervisningsskemaet, som har skubbet til hvornår vi har holdt frokostpause i den ene halvdel af året, men ikke den anden, hvilket derved udvisker de klare søjler. Det burde dog betyde at den ene høje søjle skulle være flere små søjler, men tirsdagens søjler er faktisk samme højde som mandagens veldefinerede søjler. Her skal vi dog husk på at mandag har færre pings end tirsdag, og derfor er dens søjler naturligt mindre.

% \begin{figure}[h!]
% 	\centering
% 	\resizebox{\textwidth}{!}{\input{plot_weekday_by_hour_analysis_2024-2025.pgf}}
% 	\vspace{-25pt}
% 	\caption{\protect\coffee\ pings fordelt på klokkeslet og ugedag i 2024-2025 perioden.}
% 	\label{fig:weekday_hour_analysis_2024-2025}
% \end{figure}

Men nu tænker du jo nok som den opmærksomme læser du er, at alle uger i løbet af året umuligt kan være ens, der er jo både ferier og helligdage og alt der ind i mellem, og der har du naturligvis helt ret.
Vi kan derfor tage et kig på fordelingen af \coffee\ pings over de forskellige måneder på året. Denne fordelingen kan ses på Figur~\ref{fig:month_analysis_multi}, hvor de tre år er vist side om side.
Her kan det ses at vi nogenlunde følger fordelingen i det sidste år som de andre år. Dog ses det at der er flere \coffee\ pings i oktober, november og december i år end der var sidste år.
Jeg har forsøgt at finde ud af om dette overlapper med udlandsopholdet for nogle i chatten, dog uden held. Jeg kan derfor ikke umiddelbart forklare hvorfor oktober og november ser anderledes ud. Dog kan jeg tage stigningen i december på mine egne skuldre, da jeg tydeligt husker en paper deadline vi sad med i den tid som krævede mange ture ned efter mere \coffee.
Det sagt, så har du jo nok, kære læser, lagt mærke til at den største afvigelse er for juli måned, som viser at vi år for år er blevet værre til at holde sommerferie.
Dette er kulmineret i den nyeste periode, hvor vi har sat månedelig rekord for antal \coffee\ pings på en enkelt måned.
For at forklare hvorfor dette skete behøver jeg ikke nogle antagelser, men kan direkte forklare det ved at det var i juli måned jeg havde deadline på min Ph.D. afhandlingen, og måneden efter Lasse havde sin deadline, og vi var derfor, specielt i juli, \emph{rigtigt} meget i kaffestuen.

% \begin{figure}[h!]
% 	\centering
% 	\resizebox{\textwidth}{!}{\input{plot_month_analysis_2022-2025.pgf}}
% 	\vspace{-25pt}
% 	\caption{\protect\coffee\ pings fordelt på måned i 2022-2025 perioden.}
% 	\label{fig:month_analysis_2022-2025}
% \end{figure}

% \begin{figure}[h!]
% 	\centering
% 	\resizebox{\textwidth}{!}{\input{plot_month_analysis_2024-2025.pgf}}
% 	\vspace{-25pt}
% 	\caption{\protect\coffee\ pings fordelt på måned i 2024-2025 perioden.}
% 	\label{fig:month_analysis_2024-2025}
% \end{figure}

\begin{figure}[h!]
	\centering
	\resizebox{\textwidth}{!}{\input{plot_month_analysis_multi_2022-2023_2023-2024_2024-2025.pgf}}
	\vspace{-25pt}
	\caption{\protect\coffee\ pings fordelt på måned i alle perioder. Her er den første kolonne for hver ugedag den første periode som var i 2022-2023, og så fremdeles for de to næste kolonner.}
	\label{fig:month_analysis_multi}
\end{figure}

Men nu, kære læser, er det blevet tid til at kigge på det helt store billede, nemlig antal \coffee\ pings fordelt på hver dag i løbet af året. Denne fordeling kan ses på Figur~\ref{fig:year_analysis_2024-2025}.
Her er det tydeligt at se et hul i oktober, som paser med efterårsferien. Den har dog stadig nogle enkelte \coffee\ pings, hvilket skyldes at ugedagene rykker sig en smule fra datoerne fra år til år, dog kan de ikke rykke sig med særligt meget, og uge 42 ligger derved nogenlunde, men ikke helt, på de samme datoer hvert år.
Det samme kan ikke siges om Jesus evne til at dø nogenlunde konsistent, og det er derfor næsten umuligt at afgøre på plottet hvornår påsken har ligget i løbet af årene. Dog kan det ses at der er lidt færre pings i april end i resten af forårets måneder, hvilket passer med hvor påsken oftest falder.
Hvad han dog mangler i konsistens i hans død bliver dog gjort op for hvor konsistent han bliver født, da man kan se et markant dyk i antal \coffee\ pings til juletid.
I sidste udgave blev det beskrevet at der var op til flere datoer uden nogle pings i løbet af de to år, men dette er nu blevet reduceret kraftigt, og kun juleaften og første juledag har nu ikke nogle \coffee\ pings. 
Hvis jeg dog kender fejlkilderne korrekt gælder dette kun frem til næste jul, hvor de med alt sandsynlighed vil gøre en aktiv indsats for at nå et smut forbi kaffetuen til juletid.

\begin{figure}[h!]
	\centering
	\resizebox{\textwidth}{!}{\input{plot_year_analysis_2022-2025.pgf}}
	\vspace{-25pt}
	\caption{\protect\coffee\ pings fordelt på årets dag og måned over alle tre år.}
	\label{fig:year_analysis_2024-2025}
\end{figure}

Det er i tidligere udgaver blevet nævnt at man kan se mange naturlige huller i løbet af året, når der kun kigges på data fra et enkelt år af gangen. Disse huller passer naturligt med de weekender som ligger spredt over året. 
Jeg har taget et kig på fordelingen af \coffee\ pings på årets dage for det sidste år\footnote{Dette plot er tilsvarende udeladt i denne artikel, da det både fylder en masse, men også for det meste er uinterasant.}, hvor dette stadig er gældende for langt det meste af året. Dog er der en undtagelse i juli og august måned, hvor det lader til at at der ikke findes weekender.
Der var dog weekender i disse måneder, og mangln forklare derfor simpelt ved at Lasse og jeg af naturlige årsager ikke holdt specielt meget weekend i disse måneder. Vi slet og ret stoppede med at bekymre os om hvad datoen nu engang sagde om ugedagen, men udelukkede hvad datoen sagde om tiden tilbage før deadlinen.
Meget poetisk stemmer det, i hvert fald for juli måned, meget smukt overens med den danske Ph.D. lovgivning\footnote{https://www.retsinformation.dk/eli/lta/2025/1124}, som i kapitel 7 \S 18 får beskrevet at juli ikke er en rigtig måned, og det kan vi begge blot nikke meget enig med på.

% \begin{figure}[h!]
% 	\centering
% 	\resizebox{\textwidth}{!}{\input{plot_year_analysis_2024-2025.pgf}}
% 	\vspace{-25pt}
% 	\caption{\protect\coffee\ pings fordelt på årets dag og måned i 2024-2025 perioden.}
% 	\label{fig:year_analysis_2024-2025}
% \end{figure}

Tidligere så vi at vi over de sidste tre år gennemsniteligt har lavet 8,12 \coffee\ pings om dagen. Men som vi også har set i den forrige analyse, så er der nu nogen variation på hvor mange \coffee\ pings der er hver dag. Heldigvis har statistikerne givet os nogle værktøjer til at måle bedre end blot et gennemsnit! Her kan vi først finde frem til at medianen af antal \coffee\ pings per dag er 8. Okay, måske vi kan finde bedre værktøjer.
Sidst kiggede vi på at behandle antal \coffee\ ping per dag som en sandsynlighedsfordeling, ved at tælle hvor mange procent af dagene der har hvor mange pings. Sidst valgte vi at tælle dage med 0 \coffee\ pings med, hvilket for de sidste tre år giver en fordeling med varians 41,99 (som desværre ikke kan afrundes til det skønne tal 42 med to decimaler, øv), og standartafvigelsen bliver derved 6,48. Dette modbeviser naturligt de store tals lov, som siger det burde bliver mindre med mere data, og ikke større end sidst. Kimse kimse.
I denne fordeling er der hele 211 dage over de sidste år uden nogle \coffee\ ping, hvilket udgør 19,25\% af alle dage. Jeg vil derfor udnytte at jeg ikke er datavidenkaber, og derfor gerne må fuske lidt med dataen. Eller, ændre hvad vi måler på. For hvis vi ikke er interaserede i alle dage vi ikke er her, men i stedet vil have antal \coffee\ pings for de dage vi faktisk er her, kan vi ændre problemet til at i stedet kigge på fordelingen af \coffee\ pings over de dage der har mindst 1 \coffee\ ping. Denne fordelingen kan ses på Figur~\ref{fig:days_count_distribution}.
For det første ændrer gennemsnittet sig markant til 10,06 \coffee\ pings per dag, og medianen er 10. Og så sker der det magiske af hvis man fjerner den største outlier i ens data falder variansen! Faktisk bliver den til kun 32,52, og standartafvigelsen tilsvarende 5,7. Ganske udemærket.

\begin{figure}[h!]
	\centering
	\resizebox{\linewidth}{!}{\input{plot_days_count_distribution_ignore_zero_days.pgf}}
	\vspace{-15pt}
	\caption{Fordeling af antal \protect\coffee\ pings per dag, givet at der er mindst 1 \protect\coffee\ ping, med procent antal dage over de tre år i grå og normalfordelingen i sort.}
	\label{fig:days_count_distribution}
\end{figure}

Men nu handler \coffee\ ping chatten ikke om hvornår eller hvor ofte man går i kaffetuen, men at man går i kaffestuen sammen.
Men hvordan kan vi se udelukkende på hvornår folk \coffee\ pinger om de er gået i kaffestuen sammen? Det ærlige svar er, at det kan vi ikke. Men vi kan estimere det!
Her bruger vi den bedste parameter vi har at arbejde med til vores estimat, og definerer at to \coffee\ pings $p_1$ og $p_2$, for $p_1$ pinget før $p_2$, er gået i kaffestuen sammen, hvis er der maximalt $\Delta t$ tid imellem dem. Hvis der så er yderligere maximalt $\Delta t$ tid mellem $p_2$ og næste ping $p_3$ gruperes alle tre, og så fremledes. Her er det i første udgave blevet fastlagt at
\[ \Delta t = 5 \; \text{min} \; , \]
da det virkede som den mest fornuftige parameter. Og denne parameter er vi tvunget til at fastholde, da det er blevet kutyme at man venter præcis i denne tid med at gå igen fra kaffestuen for at se om der kommer nogle \coffee\ pings man skal vente på.
Tidligere har fordelingen af grupper dog vist at der er mange flere enkelt grupper end man kunne forvente, og dette har desværre ikke ændret sig, da der for de sidste tre år i alt har været 5785 grupper i kaffestuen, hvoraf kun 2078 grupper indeholder strengt mere end én person, hvilket efterlader hele 3707 grupper som værende enkelt personers grupper, altså er næsten 2/3 af alle grupper enkelt personer.
Kigger man på fordelingen af hvor mange af hver gruppe størelse der er, falder denne også drastisk som gruppestørrelsen stiger, ja nærmest med eksponentiel aftagning!
Dette er dog utroligt misvisende, da en gruppe på to personer jo ikke blot er en gruppe i kaffestuen, men faktisk er to personer i kaffetuen. Hvis man derfor opjusterer sit plot ved at skalere til hvor mange \coffee\ pings der i alt hører til i hver gruppe, falder den nærmere kun lineart som antal personer i gruppen stiger.

Man kan dog få nogle interasante opservationer ved at kigge på hvordan disse grupper er fordelt over dagen. Her definerer vi at tidpunktet for en gruppe er når det første \coffee\ ping i gruppen sendes. Da de fleste grupper er med 1, 2 eller 3 personer, hvor kun 208 grupper er på 4 eller flere personer, vil vi bunke 3 eller flere personer sammen. Hvis vi blot kigger på antal grupper vil grupper med eks 2 personer automatisk være mindre end 1 personer, da det kræver flere \coffee\ pings at lave en 2 personers gruppe, tilsvarend ti argumntet fra før. Vi vil derfor justere fordelingen til at ikke vise antallet af grupper, men antallet af personer eller \coffee\ pings der er i hver bunke. Dette plot kan ses på Figur~\ref{fig:group_size_by_hour}.
Bemærk at plottet er skåret til en udvidet arbejdsdag fra kl 7 til 18. Her er der tre serier, som for hver time viser antal personer i gruppernes fordeling.
Vi kan se at om morgenen er der mange enkelt personers grupper, hvilket kan forklares ved at mange lige skal omkring kaffestun når de møder ind, men at de ikke møder ind på samme tid, og derved ikke går i kaffestuen sammen. Derimod holder vi tit fælles frokost pause, hvilket kan ses på at der kl 12 og 13 er langt flere personer der følges ad i grupper.
Dataen, der er undladt, viser at der er meget få pings før kl 7, og de få der er er primært 1 personers grupper, mens fordelingen efter kl 18 ligner kl 18 fordelingen, dog uden så mange 3+ personers grupper.

\begin{figure}[h!]
	\centering
	\resizebox{\textwidth}{!}{\input{plot_groups_per_hour.pgf}}
	\vspace{-25pt}
	\caption{\protect\coffee\ pings fordelt på gruppestørrelser og timer i løbet af dagen. Den venstre serie er 1 personers grupper, den midterste 2 personer, og den højre 3 eller flere personer.}
	\label{fig:group_size_by_hour}
\end{figure}


Sidst kiggede vi på et kumulativ plot over antal \coffee\ pings over tiden der er gået, og bemærkede at der var et fald i kurven da Steffan var afsted på sit udlandsophold. Hertil blev det naturligvis konkluderet at faldet var Steffans skyld. Dog kunne vi se at der manglede ca. 1000 \coffee\ pings, mens Steffan i samme periode kun burde have produceret ca. 200 \coffee\ pings.
Konklusionen måtte derved være at Steffan trak ca. 4 personer med sig i kaffestuen når han pingede, dog blev det aldrig målt.
Men hvordan skulle man dog måle det? Ja, her må man tage en side ud af den lille røde bog, og implementere et social credit system i \coffee\ ping chatten. Til dette skal vi have tillagt os et mål for hvad der er god og dårlig opførsel, samt hvor god og dårlig den er, før vi kan lave et sådan system. Det er naturligvis godt at gå i kaffestuen. Derudover er det godt at trække andre mennesker med sig i kaffestuen. Her kan vi bruge gruppe analysen fra før til at bestemme hvem man trækker med sig, ved at definere at alle i gruppen der kommer efter en, har man trukket med. Set fra det omvendte perspektiv, så bliver man trukket af alle personer der kommer før en i gruppen. Dog kan de vel ikke alle trække lige meget, da den første person i gruppen naturligt er lederen som trækker mest, og dem er man trukket af. Så er der endnu en der pinger og er den næste i gruppen, og dem er man også trukket af, men da de allerede er trukket af lederen, trækker de ikke i samme omfang på en.
Hvis vi tager en side ud af matematikernes bog, kan vi med fordel bruge en geometrisk aftagende fordeling, som basis for hvor meget man trækker. Ved at definere første person til at trække i en med kraft 1 og den næste med kraft $\frac{1}{2}$, den næste med kraft $\frac{1}{4}$, og den $i$'te (naturligvis 0-indekseret) trækker med $\frac{1}{2^i}$, så ved vi at man bliver trukket i med maximalt kraft 2, og derved bliver den totale trækkerkraft maximalt 2 gange antal \coffee\ pings. Ved at kigge tilbage på den sociale credit score, kan vi definere at hvor meget man totalt trækker i andre, er et godt mål for ens score.
Ved at køre denne analyse, får man først en rå score ud, som for Steffan er ca. 622,52. Dette betyder dog ikke så meget for diskutionen om hvor meget han trækker per gang han \coffee\ pinger. Derfor kan man normalisere scoren med antal grupper man er med i, for at få et et mål for trækkraften per gruppe. Her kommer Steffans normaliserede score frem til ca. 0,45. Han trækker derved en halv person med sig når han går i kaffestuen, altså en del mindre end de formodede 4 personer.
Hvis vi ændrer vores trækkraft til at være ens for alle tidligere personer, med den naturlige idé om at flere mennesker trækker mere, og derved alle kan trække med 1\footnote{Undskyld matematikere for at fjerne den pæne geometriske aftagning. Dog læg lige mærke til at denne fodnote kunne være en potens, og sætningen ville stadig være ens, magisk.}, bliver Steffans normaliserde score til ca. 0,53. Pokkers, vores mål virker ikke til at forklare de manglende pings.

Man kan derfor vende sig mod et eftertestet mål, brugt over hele verden, nemlig Elo score. Her har hver prson en score, og når de går i kamp mod hinanden beregner Elo scoren den ændring som scorene skal have, som en funktion af hvor sandsynligt det er for hver at vinde. Dvs, hvis man vinder til nogen med en lavere score, så var det sandsynligt det skete, og derfor stiger ens score ikke meget, men hvis man vinder til nogen med en meget højere score, får man også en højere ændring.
Vi skal derfor først definere hvad det vil sige at lave en kamp, og hvad det vil sige at vinde eller tabe den kamp. Her har jeg valgt at alle i en gruppe spiller en kamp mod hinanden, hvor vinderen er den som pinger først i gruppen.
Her er start scoren 1200 for hver person, og vi sætter sværhedsgradsparameteren $K$ til 16, som er det Wikipedia mener den skal være for eksperter, da vi med vores antal \coffee\ ping så sandeligt er eksperter. Efter alle kampene hen over de sidste 3 år ender Steffan med en Elo score på -7740,13, og er derved den person med den suverent mindste score, mens topscoren har en Elo score på 8476,87. Ugf, det er jo et endnu værre resultat end målingen med social credit score!
Men! Nu skal vi huske, at der i en Elo kamp både er en vinder og en taber, og taberen mister point, hvilket også er hvorfor Steffan ender så lavt. Men alle kampe bunder altid ud i at vi går i kaffestuen, så kan man dog overhovedet tale om at man kan tabe en kamp? Nej da! Man kan kun vinde!
Ved at lade scoren starte på 0, og ikke sænke scoren hvis man taber en kamp, vender billedet sig fuldstændigt, og Steffan ender med den suverent største score på 9976,84. Vanvittigt! Til sammenligning er den næsthøjeste score på 7395,37. Steffan er derfor med dette mål den suverent vildeste trækker af folk i kaffestuen, og vi kan derfor konkludere, at det er det bedste mål. Hvad tallet har af indvirkning på virkeligheden er dog uklart, men det er en mindre detalje.


\section*{Konklusion}

Ovenpå enhver god analyse følger der naturligt en konklusion. Men hvad kan man overhovedet konkludere ovenpå et sådan virvar af plots og tal? Kan man overhovedet opsumere alt vi har kigget på til en enkelt linjes pæn tekst? Hvis man kan, kender jeg den i hvert fald ikke, andet end at vi muligvis går for meget i kaffestuen.
Jeg håber derfor på at du kære læser selv har nogle konklussioner, du vil sidde tilbage med efter at have kigget mine tanker og plots igennem, for ellers virker det måske lidt som spild af tid at have læst.
Det sagt; skal vi gå i kaffestuen? \coffee\

\end{document}
